\uuid{tr6h}
\exo7id{4407}
\titre{exo7 4407}
\auteur{quercia}
\organisation{exo7}
\datecreate{2010-03-12}
\isIndication{false}
\isCorrection{false}
\chapitre{Série numérique}
\sousChapitre{Fonction exponentielle complexe}
\cpge{mpsi-mp2i.fonctions-valeurs-reelles-complexes.d, pcsi.techniques-calcul-fonctions.d, ptsi.techniques-calcul-fonctions.d, mpsi-mp2i.primitives-equations-differentielles.a, pcsi.techniques-calcul-primitives.a, ptsi.techniques-calcul-primitives.a, mpsi-mp2i.derivabilite.d, pcsi.derivabilite.d, ptsi.derivabilite.d, bcpst1.derivation.a, bcpst1.derivation.c}
\module{Analyse}
\niveau{L1}
\difficulte{}

\contenu{
\texte{
Soit $f : \R \to {\C^*}$ telle que :
$\forall\ x,y\in\R,\ f(x+y) = f(x)f(y)$.
}
\begin{enumerate}
    \item \question{Si $f$ est dérivable, montrer qu'il existe $\lambda \in \C$ tel que :
    $\forall\ x \in \R,\ f(x) = e^{\lambda x}$.}
    \item \question{Obtenir le même résultat si $f$ est seulement supposée continue
    (prendre une primitive, $F$, de $f$ et montrer qu'elle est de classe $\C^2$).}
\end{enumerate}
}
